% Part: many-valued-logic
% Chapter: three-valued-logics
% Section: lukasiewicz

\documentclass[../../../include/open-logic-section]{subfiles}

\begin{document}

\olfileid{mvl}{thr}{luk}

\olsection{\L ukasiewicz-এর যুক্তিবিদ্যা}

বহুমানী যুক্তিবিদ্যার প্রথম দিকের প্রকাশিত ও বিশদভাবে বিকশিত প্রস্তাবগুলির
একটি ১৯২১ সালে পোলীয় দার্শনিক ইয়ান \L ukasiewicz দেন। অ্যারিস্টটলের
সমুদ্রযুদ্ধের সমস্যা তাঁকে অনুপ্রাণিত করেছিল: মনে হয়,
\emph{আজ} “আগামীকাল একটি সমুদ্রযুদ্ধ হবে” বাক্যটি সত্যও নয়, মিথ্যাও
নয়—তার সত্যমান এখনও স্থির হয়নি। \L ukasiewicz এই ধরনের
“ভবিষ্যৎ-আপতিক” বাক্যের জন্য তৃতীয় একটি অনির্ধারিত সত্যমান প্রবর্তনের
প্রস্তাব দেন।
\begin{quote}
আমি স্ববিরোধ ছাড়াই ধরে নিতে পারি যে আগামী বছরের কোনো নির্দিষ্ট মুহূর্তে,
যেমন ২১ ডিসেম্বর দুপুরে, ওয়ারশতে আমার উপস্থিতি বর্তমানে ইতিবাচকভাবেও
নির্ধারিত নয়, নেতিবাচকভাবেও নয়। অতএব নির্দিষ্ট সময়ে আমার ওয়ারশতে
উপস্থিত থাকা সম্ভব, কিন্তু আবশ্যিক নয়। এই অনুমানের ভিত্তিতে “আগামী বছরের
২১ ডিসেম্বর দুপুরে আমি ওয়ারশতে থাকব” বচনটি বর্তমানে সত্যও হতে পারে না,
মিথ্যাও হতে পারে না। কারণ এটি এখন সত্য হলে ওয়ারশতে আমার ভবিষ্যৎ উপস্থিতি
আবশ্যিক হতে হতো, যা অনুমানের বিরোধী। অন্যদিকে, এটি এখন মিথ্যা হলে ওয়ারশতে
আমার ভবিষ্যৎ উপস্থিতি অসম্ভব হতে হতো, যা অনুমানেরও বিরোধী। তাই আলোচ্য
বচনটি এই মুহূর্তে সত্যও নয়, মিথ্যাও নয়; এবং “0” বা মিথ্যাত্ব ও~“1”
বা সত্যতা থেকে পৃথক তৃতীয় একটি মান তার থাকতেই হবে। এই মানকে আমরা
“$\frac{1}{2}$” দিয়ে চিহ্নিত করতে পারি। এটি “সম্ভব”-কে উপস্থাপন করে এবং
তৃতীয় মান হিসেবে “সত্য” ও~“মিথ্যা”-র সঙ্গে যুক্ত হয়।
\end{quote}
\L ukasiewicz-এর তৃতীয় সত্যমানের জন্য আমরা~$\Undef$ ব্যবহার করব।%
\footnote{এখানে \L ukasiewicz “সম্ভব” শব্দটি আজকাল অপ্রচলিত এক অর্থে,
অর্থাৎ সম্ভব কিন্তু আবশ্যিক নয়—এই অর্থে ব্যবহার করেছেন।}

এই ব্যাখ্যায় $\lnot$, $\land$ ও~$\lor$ সংযোজকের সত্যমান-অপেক্ষকগুলি
সহজেই নির্ধারণ করা যায়: কোনো ভবিষ্যৎ-আপতিক বাক্যের নঞর্থক রূপটিও
ভবিষ্যৎ-আপতিক বাক্য; তাই $\tf{\lnot}(\Undef) = \Undef$। কোনো
সংযোজনের একটি সংযোজ্য অনির্ধারিত এবং অন্যটি সত্য হলে সংযোজনটিও
অনির্ধারিত—কারণ ভবিষ্যৎ-আপতিক সংযোজ্যটি শেষ পর্যন্ত কী হয় তার উপর নির্ভর
করে সংযোজনটি সত্যও হতে পারে, মিথ্যাও হতে পারে। তাই
\[
    \tf{\land}(\True, \Undef) =
\tf{\land}(\Undef, \True) =
\Undef.
\]
কিন্তু অন্য সংযোজ্যটি মিথ্যা হলে সংযোজনটি শেষ পর্যন্ত সত্য হতে পারে না;
তাই
\[\tf{\land}(\False, \Undef) =
\tf{\land}(\Undef, \False) = \False.\]
অন্য মানগুলি—অর্থাৎ আর্গুমেন্টগুলি যখন নির্ধারিত সত্যমান $\True$ অথবা
$\False$—ধ্রুপদি যুক্তিবিদ্যার মতোই।

শর্তবচনের ক্ষেত্রে পরিস্থিতি একটু জটিল। ধরা যাক, $q$ একটি
ভবিষ্যৎ-আপতিক উক্তি। $p$ মিথ্যা হলে $q$ শেষ পর্যন্ত যাই হোক না কেন,
$p \lif q$ সত্য হবে; তাই আমাদের $\tf{\lif}(\False, \Undef) = \True$
স্থির করা উচিত। আবার $p$ সত্য হলে $q$ শেষ পর্যন্ত যাই হোক না কেন,
$q \lif p$ সত্য হবে; তাই $\tf{\lif}(\Undef, \True) = \True$। $p$ সত্য
হলে $p \lif q$ শেষ পর্যন্ত সত্য বা মিথ্যা—দুই-ই হতে পারে; তাই
$\tf{\lif}(\True, \Undef) = \Undef$। একইভাবে, $p$ মিথ্যা হলে
$q \lif p$ শেষ পর্যন্ত সত্য বা মিথ্যা হতে পারে; তাই
$\tf{\lif}(\Undef, \False) = \Undef$। এখন কেবল $p$ ও~$q$ উভয়ই
ভবিষ্যৎ-আপতিক হওয়ার ক্ষেত্রটি বাকি। প্রেরণাটিকে অনুসরণ করলে এখানে
আসলে~$\Undef$ দেওয়া উচিত। কিন্তু তাতে $!A \lif !A$
\emph{সর্বতঃসত্য হবে না}। \L ukasiewicz-এর $!A \lor \lnot !A$ এবং
$\lnot(!A \land \lnot !A)$ ত্যাগ করতে আপত্তি ছিল না, কিন্তু
$!A \lif !A$ ত্যাগ করতে তিনি রাজি হননি। তাই তিনি
$\tf{\lif}(\Undef, \Undef) = \True$ স্থির করেন।

\begin{defn}\ollabel{def:lukasiewicz}
ত্রিমানী \L ukasiewicz যুক্তিবিদ্যা নিচের ম্যাট্রিক্স দিয়ে সংজ্ঞায়িত:
\begin{enumerate}
  \item $\lnot$, $\land$, $\lor$, $\lif$-সহ প্রমিত বচনগত ভাষা
  $\Lang L_0$।
  \item সত্যমানের সেট $V = \{\True, \Undef, \False\}$।
  \item $\True$ একমাত্র মনোনীত মান; অর্থাৎ $V^+ = \{\True\}$।
  \item সত্যমান-অপেক্ষকগুলি নিচের সারণিগুলি দিয়ে দেওয়া হয়:
  \begin{center}
    \begin{tabular}{c|c}
      $\tf{\lnot}$ & \\
      \hline
      $\True$ & $\False$ \\
      $\Undef$ & $\Undef$ \\
      $\False$ & $\True$
    \end{tabular}
    \quad
    \begin{tabular}{c|ccc}
      $\tf{\land}[\LogLuk[3]]$ & $\True$ & $\Undef$ & $\False$ \\
      \hline
      $\True$ & $\True$ & $\Undef$ & $\False$ \\
      $\Undef$ & $\Undef$ & $\Undef$ & $\False$\\
      $\False$ & $\False$ & $\False$ & $\False$
    \end{tabular}
    \\[2ex]
    \begin{tabular}{c|ccc}
      $\tf{\lor}[\LogLuk[3]]$ & $\True$ & $\Undef$ & $\False$ \\
      \hline
      $\True$ & $\True$ & $\True$ & $\True$ \\
      $\Undef$ & $\True$ & $\Undef$ & $\Undef$ \\
      $\False$ & $\True$ & $\Undef$ & $\False$
    \end{tabular}
    \quad
    \begin{tabular}{c|ccc}
      $\tf{\lif}[\LogLuk[3]]$ & $\True$ & $\Undef$ & $\False$ \\
      \hline
      $\True$ & $\True$ & $\Undef$ & $\False$ \\
      $\Undef$ & $\True$ & $\True$ & $\Undef$  \\
      $\False$ & $\True$ & $\True$ & $\True$
    \end{tabular}
  \end{center}
\end{enumerate}
\end{defn}

সহজেই দেখা যায়, কেবল $\lnot$, $\land$ ও~$\lor$ থাকা যেকোনো
!!{formula}~$!A$-এর সব !!{propositional variable}s-কে~$\Undef$ আরোপ
করলে সূত্রটির সত্যমানও~$\Undef$ হবে। তাই, যেমন, ধ্রুপদি সর্বতঃসত্য
$p \lor \lnot p$ ও~$\lnot(p \land \lnot p)$, $\LogLuk[3]$-এ সর্বতঃসত্য
নয়; কারণ $\pAssign v(p) = \Undef$ হলে উভয় ক্ষেত্রেই
$\pValue{v}(!A) = \Undef$।

যেসব !!{valuation}s-এ $\pAssign v(p) = \True$ অথবা~$\False$, সেখানে
$\pValue v(!A)$ তার ধ্রুপদি সত্যমানের সঙ্গে মিলে যাবে।

\begin{prop}
  $!A$-তে থাকা প্রতিটি $p$-এর জন্য
  $\pAssign v(p) \in \{\True, \False\}$ হলে
  $\pValue v(!A)[\LogLuk[3]] = \pValue v(!A)[\LogCL]$।
\end{prop}

\begin{prob}\label{mvl:thr:luk:prob:luk-iff} ধরা যাক, $\LogLuk[3]$-এ
  আমরা $\pValue v(!A \liff !B) = \pValue v((!A \lif !B) \land (!B \lif
  !A))$ সংজ্ঞায়িত করি। তাহলে $\liff$-এর সত্যসারণি কী হবে?
\end{prob}

অনেক ধ্রুপদি সর্বতঃসত্য \LogLuk[3]-এও সর্বতঃসত্য; যেমন,
$\lnot p \lif (p \lif q)$। ধ্রুপদি যুক্তিবিদ্যার মতোই আমরা সত্যসারণি
দিয়ে এটি যাচাই করতে পারি:
\begin{center}
\begin{tabular}{cc|cccccc}
  $p$ & $q$ & $\lnot$ & $p$ & $\lif$ & $\smash{(}p$ & $\lif$ & $q\smash{)}$ \\ \hline
  \True & \True & \False & \True & \True & \True & \True & \True \\
  \True & \Undef & \False & \True & \True & \True & \Undef & \Undef \\
  \True & \False & \False & \True & \True & \True & \False & \False \\
  \Undef & \True & \Undef & \Undef & \True & \Undef & \True & \True \\
  \Undef & \Undef & \Undef & \Undef & \True & \Undef & \True & \Undef \\
  \Undef & \False & \Undef & \Undef & \True & \Undef & \Undef & \False \\
  \False & \True & \True & \False & \True & \False & \True & \True \\
  \False & \Undef & \True & \False & \True & \False & \True & \Undef \\
  \False & \False & \True & \False & \True & \False & \True & \False \\
\end{tabular}
\end{center}

\begin{prob}
  দেখাও যে নিচের সূত্রগুলি \LogLuk[3]-এ সর্বতঃসত্য:
  \begin{enumerate}
    \item $p \lif (q \lif p)$
    \item\label{mvl:thr:luk:prob:luk-taut-2} $\lnot(p \land q) \liff (\lnot p \lor \lnot q)$
    \item\label{mvl:thr:luk:prob:luk-taut-3} $\lnot(p \lor q) \liff (\lnot p \land \lnot q)$
  \end{enumerate}
  (\olref[mvl][thr][luk]{prob:luk-taut-2} ও
  \olref[mvl][thr][luk]{prob:luk-taut-3}-এ $!A \liff !B$-কে
  $(!A \lif !B) \land (!B \lif !A)$-এর সংক্ষেপ হিসেবে নাও, অথবা
  \olref[mvl][thr][luk]{prob:luk-iff}-এর তোমার সমাধান ব্যবহার করো।)
\end{prob}

\begin{prob}
  দেখাও যে নিচের ধ্রুপদি সর্বতঃসত্যগুলি~\LogLuk[3]-এ সর্বতঃসত্য নয়:
  \begin{enumerate}
    \item $(\lnot p \land p) \lif q$
    \item $((p \lif q) \lif p) \lif p$
    \item $(p \lif (p \lif q)) \lif (p \lif q)$
  \end{enumerate}
\end{prob}

তাই কেউ হয়তো ভাবতে পারেন যে সব ধ্রুপদি সর্বতঃসত্য~$\LogLuk[3]$-এ
সর্বতঃসত্য না হলেও, প্রতিটি !!{valuation}-এ অন্তত তাদের মান~$\True$ অথবা
$\Undef$ হবে। তা নয়। একটি প্রতিদৃষ্টান্ত হলো
\[
  \lnot(p \lif \lnot p) \lor \lnot(\lnot p \lif p)
\]
যার মান $p$-এর মান~$\Undef$ হলে~$\False$।

\begin{prob}
  \L ukasiewicz যুক্তিবিদ্যায় নিচের কোন সম্পর্কগুলি সিদ্ধ? প্রতিটির
  জন্য একটি সত্যসারণি দাও।
  \begin{enumerate}
    \item $p, p \lif q \Entails q$
    \item $\lnot\lnot p \Entails p$
    \item $p \land q \Entails p$
    \item $p \Entails p \land p$
    \item $p \Entails p \lor q$
  \end{enumerate}
\end{prob}

\L ukasiewicz তাঁর ত্রিমানী ব্যবস্থার ভিত্তিতে সম্ভাব্যতার একটি যুক্তিবিদ্যা
গড়তে চেয়েছিলেন। এজন্য তিনি এক-স্থানীয় সংযোজক $\Diamond !A$
(“$!A$ সম্ভব”) এবং তার অনুরূপ $\Box !A$ (“$!A$ আবশ্যিক”) প্রবর্তন করেন:
\begin{center}
  \begin{tabular}{c|c}
    $\tf{\Diamond}$ & \\
    \hline
    $\True$ & $\True$ \\
    $\Undef$ & $\True$ \\
    $\False$ & $\False$
  \end{tabular}
  \quad
  \begin{tabular}{c|c}
    $\tf{\Box}$ & \\
    \hline
    $\True$ & $\True$ \\
    $\Undef$ & $\False$ \\
    $\False$ & $\False$
  \end{tabular}
\end{center}
অন্যভাবে বললে, $p$ ঠিক তখনই সম্ভব, যখন সেটি ইতিমধ্যেই মিথ্যা বলে স্থির
হয়নি; আর $p$ ঠিক তখনই আবশ্যিক, যখন সেটি ইতিমধ্যেই সত্য বলে স্থির হয়েছে।

\begin{prob}
  দেখাও যে $\Box p \liff \lnot\Diamond \lnot p$ এবং $\Diamond p \liff
  \lnot \Box \lnot p$, $\Box$ ও~$\Diamond$-এর সত্যসারণি দিয়ে প্রসারিত
  \LogLuk[3]-এ সর্বতঃসত্য।
\end{prob}

কিন্তু প্রস্তাবিত এই মোডাল যুক্তিবিদ্যার সীমাবদ্ধতা শিগগিরই স্পষ্ট হয়:
ঘটনা যেভাবেই ঘটুক, $p \land \lnot p$ কখনও সত্য হতে পারে না। তাই এখন
এটি স্থির না হলেও (এবং সেই কারণে অনির্ধারিত হলেও) একে অসম্ভব গণ্য করা
উচিত; অর্থাৎ $\lnot \Diamond(p \land \lnot p)$ সর্বতঃসত্য হওয়া উচিত।
কিন্তু $\pAssign v(p) = \Undef$ হলে
$\pValue v(\lnot \Diamond(p \land \lnot p)) = \False$। সম্ভাব্যতা ও
আবশ্যিকতার মতো মোডাল পার্থক্য ধারণ করার জন্য দুটি সত্যমান যথেষ্ট নয়—এই
বিষয়ে \L ukasiewicz ঠিক ছিলেন; কিন্তু তৃতীয় একটি সত্যমান যোগ করাও
যথেষ্ট নয়।
% BN-SRC-318 TARGET-CORRECTION-BEGIN
\par\smallskip\noindent\textnormal{\small\textbf{উৎস-সংশোধন (BN-SRC-318):} হিমায়িত উৎসে মিথ্যা ও অনির্ধারিত সংযোজ্যের সমমিত ক্ষেত্র দেখাতে $\tf{\land}(\False,\Undef)$ একই আর্গুমেন্ট-ক্রমে দুবার লেখা হয়েছে। সত্যসারণি ও বাক্যের অভিপ্রায় অনুযায়ী দ্বিতীয়টিকে $\tf{\land}(\Undef,\False)$ হিসেবে বাংলা লক্ষ্যে পুনরুদ্ধার করা হয়েছে।} % BN-SRC-318 TARGET-CORRECTION-END
% BN-SRC-319 TARGET-CORRECTION-BEGIN
\par\smallskip\noindent\textnormal{\small\textbf{উৎস-সংশোধন (BN-SRC-319):} হিমায়িত অনুশীলনে সূত্র $\bigl((\lnot p\land p)\lif q\bigr)$-এর শেষে উদ্বোধনী জোড়াবিহীন একটি অতিরিক্ত ডান বন্ধনী ছিল। বাংলা লক্ষ্যে অতিরিক্ত বন্ধনীটি বাদ দিয়ে সুষম সূত্রটি পুনরুদ্ধার করা হয়েছে।} % BN-SRC-319 TARGET-CORRECTION-END
% BN-SRC-323 TARGET-CORRECTION-BEGIN
\par\smallskip\noindent\textnormal{\small\textbf{উৎস-সংশোধন (BN-SRC-323):} হিমায়িত উৎসে $\pAssign v(p)=\Undef$ হলে $\pValue v(\lnot\Diamond(p\land\lnot p))=\Undef$ বলা হয়েছে। প্রদত্ত সারণি অনুসারে $\lnot p=\Undef$, $p\land\lnot p=\Undef$, $\Diamond\Undef=\True$, এবং $\lnot\True=\False$; তাই বাংলা লক্ষ্যে শেষ মানটি~$\False$ করা হয়েছে। এতে সূত্রটি সর্বতঃসত্য নয়—এই অভীষ্ট সিদ্ধান্তই বজায় থাকে।} % BN-SRC-323 TARGET-CORRECTION-END

\end{document}
